\section{Hardware y materiales}
\label{sec:hardware}

El robot está armado sobre un chasis diferencial de dos ruedas motrices (2WD) controlado por un ESP32-S3~\cite{espressif2024esp32s3}. En esta sección detallamos tanto la configuración base con la que corrimos las primeras pruebas en \texttt{main} como los cambios de potencia y lectura de batería que añadimos después en la rama \archivo{origin/firmware/volver}~\cite{repo2026agentelaberinto,arquitectura2026doc}.

% ============================================================================
\subsection{Lista de materiales (BOM) y elección de componentes}
\label{subsec:bom}

Cuando armamos el carro tuvimos que equilibrar tres cosas: que rodara sin cables, que pudiera leer sensores analógicos con el Wi-Fi prendido a \SI{2.4}{\giga\hertz} y que cupiera con holgura dentro de celdas de \SI{24}{\centi\meter}. La \Cref{tab:bom} recoge todos los componentes que terminaron en el chasis.

\begin{table}[H]
  \centering
  \caption{Lista de materiales (BOM) y especificaciones técnicas de la plataforma física.}
  \label{tab:bom}
  \renewcommand{\arraystretch}{1.18}
  \footnotesize
  \begin{tabularx}{\textwidth}{@{}p{2.5cm} p{2.9cm} c X p{3.1cm}@{}}
    \toprule
    \textbf{Subsistema} & \textbf{Componente / Ref.} & \textbf{Cant.} & \textbf{Especificaciones técnicas clave} & \textbf{Función en el sistema} \\
    \midrule
    Cómputo y comunicaciones &
      ESP32-S3 N16R8 (44 pines) + Shield de expansión &
      1 &
      \SI{16}{\mega\byte} Flash, \SI{8}{\mega\byte} OPI PSRAM, Wi-Fi \SI{2.4}{\giga\hertz}, OTA (\cmd{app3M\_fat9M\_16MB}), USB-C UART CH343 (\cmd{DEC-06}), jack DC \SIrange{6.5}{9}{\volt} &
      Planificación a bordo, control en lazo cerrado, consola TCP y OTA \\
    \addlinespace[3pt]
    Estructura mecánica &
      Chasis 2WD + ruedas + rueda loca &
      1 &
      Plataforma $18 \times \SI{10}{\centi\meter}$; ancho con llantas \SI{15}{\centi\meter}; largo total \SI{20}{\centi\meter}; llantas $\varnothing\,\SI{43}{\milli\meter} \times \SI{19}{\milli\meter}$, eje D de \SI{3}{\milli\meter} &
      Soporte rígido de sensores, driver, IMU y packs de potencia \\
    \addlinespace[3pt]
    Actuación &
      Motorreductor N20-10 &
      2 &
      \SI{6}{\volt} nominales, \SI{500}{RPM}, sin encoders; velocidad tangencial nominal $\approx \SI{1.12}{\meter\per\second}$ &
      Tracción diferencial derecha (Motor~A) e izquierda (Motor~B) \\
    \addlinespace[3pt]
    Etapa de potencia &
      Driver L298N &
      1 &
      Doble puente H bipolar; caída interna $\approx \SI{2}{\volt}$; jumpers \cmd{ENA}/\cmd{ENB}/\cmd{5V} retirados; borne \SI{5}{\volt} desde shield &
      Control PWM de velocidad y sentido de giro de ambos motores \\
    \addlinespace[3pt]
    Percepción óptica &
      Sensor IR TCRT5000 &
      3 &
      Salida analógica (\cmd{AO}), alimentación \SI{3.3}{\volt} (filas G/V/S), ADC1 de 12~bits (\numrange{0}{4095}) &
      Detección de líneas negras (izquierdo, frontal y derecho) \\
    \addlinespace[3pt]
    Sensor inercial &
      MPU6050 (módulo GY-521) &
      1 &
      Bus I2C a \SI{400}{\kilo\hertz} (\cmd{0x68}), giroscopio eje $Z$ ($\pm\SI{500}{\degree\per\second}$, filtro \SI{44}{\hertz}), cinta doble faz de espuma &
      Integración de rumbo absoluto para giros de \SI{90}{\degree} y rectitud \\
    \addlinespace[3pt]
    Señalización &
      LED RGB WS2812 &
      1 &
      Integrado en GPIO~48 (puente de soldadura \cmd{RGB} cerrado en placa clon) &
      Indicador visual de estado, mapa, movimiento, OTA y \emph{brownout} \\
    \addlinespace[3pt]
    Alimentación lógica &
      Batería alcalina \SI{9}{\volt} + conector jack DC &
      1 &
      \SI{9.5}{\volt} nueva ($\approx \SI{8.5}{\volt}$ umbral mínimo de cambio); potencia Wi-Fi limitada a \SI{8.5}{dBm} &
      Alimentación exclusiva del ESP32-S3 y sensores (aislada de picos) \\
    \addlinespace[3pt]
    Alimentación motores &
      Pack $4\times\text{AA}$ \badgeok{} / 2S 18650 \badgepr{} &
      1 &
      $4\times\text{AA}$ (\SI{6}{\volt}) en \cmd{main}; pack 2S 18650 (\SIrange{7.8}{8.4}{\volt}) con interruptor general en \archivo{firmware/volver} &
      Alimentación del riel de potencia (\SI{12}{\volt}/VMS) del L298N \\
    \addlinespace[3pt]
    Telemetría de batería \badgepr &
      Divisor resistivo + filtro RC &
      1 &
      Resistencias \SI{6.6}{\kilo\ohm} / \SI{2.2}{\kilo\ohm} + capacitor \SI{10}{\nano\farad} en GPIO~1 (ADC1); relación calibrada $4.195$ &
      Medición de $V_{\text{bat}}$ en movimiento y compensación dinámica PWM \\
    \bottomrule
  \end{tabularx}
\end{table}

\paragraph{El problema de los motores N20 rápidos y el L298N.}
Con llantas de \SI{43}{\milli\meter}, unos N20-10 de \SI{500}{RPM} dan en vacío una velocidad lineal de:
\begin{equation}
  v_{\text{nom}} = \frac{\pi \cdot D \cdot \text{RPM}}{60} = \frac{\pi \cdot \SI{0.043}{\meter} \cdot 500}{60} \approx \SI{1.125}{\meter\per\second}.
  \label{eq:vel-nominal}
\end{equation}
Sobre el papel parecen ágiles, pero en un laberinto pequeño eso jugó en contra. Demasiada velocidad implica poco par, y además estos motores venían sin encoders~\cite{siegwart2011autonomous}. A eso se sumó el L298N: sus transistores bipolares Darlington se comen cerca de \SI{2}{\volt}~\cite{st2000l298n}, así que con cuatro pilas AA (\SI{6}{\volt}) a los motores apenas les llegan unos \SI{4}{\volt}. En los giros sobre el eje, donde las dos ruedas arrastran fricción lateral, por debajo de \cmd{pwm\_giro = 170} los motores solo pitan y no se mueven; tuvimos que subir el arranque a \cmd{230} sobre 255, mientras que para avanzar usamos \cmd{pwm = 180} y \cmd{pwm\_aprox = 140} al llegar a una pared. Como no medimos distancia en las ruedas, el avance depende de tiempo calibrado (\cmd{t\_celda}), corregido en dirección por el MPU6050 y reajustado al toparse con paredes frontales~\cite{arquitectura2026doc}.

\paragraph{Cableado del L298N y LED de estado.}
En el L298N hay que quitar los tres \emph{jumpers} (\cmd{ENA}, \cmd{ENB} y \cmd{5V}) y llevar \SI{5}{\volt} desde el shield al borne \cmd{5V} del driver. Con $4\times\text{AA}$ su regulador interno no alcanza a entregar \SI{5}{\volt}; si ese cable se suelta, la lógica del L298N se apaga y el carro deja de moverse aunque las pilas estén llenas~\cite{handoff2026doc}. El Motor~A (\cmd{OUT1--OUT2}) mueve la llanta derecha y el Motor~B (\cmd{OUT3--OUT4}) la izquierda.

Para no depender siempre de la consola, usamos el LED RGB WS2812 del GPIO~48 (en esta placa clon hubo que cerrar con estaño el puente marcado \cmd{RGB}). Azul significa que no hay mapa cargado, verde que la ruta está lista, blanco que el carro está ejecutando \cmd{GO}, rojo que el último reinicio fue por \emph{brownout} (\cmd{ESP\_RST\_BROWNOUT}) y magenta que está recibiendo una actualización OTA~\cite{repo2026agentelaberinto}.

% ============================================================================
\subsection{Geometría del robot y dimensiones del laberinto (\cmd{DEC-02})}
\label{subsec:geometria}

La \Cref{fig:geometria-carro} muestra las medidas del robot centrado en una celda estándar, tomando el punto medio del eje motriz como origen $(0,0)$.

\begin{figure}[H]
  \centering
  % ============================================================================
% Figura: Plano técnico acotado de la geometría del robot en celda de 24 cm
% ============================================================================
\begin{tikzpicture}[
  x=0.40cm, y=0.40cm,
  >=Stealth,
  cota/.style={<->, thick, azuldoc},
  ext/.style={thin, gristexto!70, densely dashed},
  linea/.style={fill=black!88},
  chasis/.style={fill=azulfondo, draw=azuldoc, line width=1.1pt, rounded corners=2pt},
  rueda/.style={fill=gristexto!85, draw=black, line width=0.8pt, rounded corners=1.5pt},
  sensor/.style={fill=rojoborde, draw=black, line width=0.6pt}
]

  % --- Fondo blanco de la celda y margen exterior ---
  \fill[white] (-16, -15) rectangle (16, 19);

  % --- Líneas negras del laberinto (celda 24 cm entre centros, línea 2 cm) ---
  % Línea izquierda: centro en x = -12 (de -13 a -11)
  \fill[linea] (-13, -13) rectangle (-11, 13);
  % Línea derecha: centro en x = +12 (de +11 a +13)
  \fill[linea] (11, -13) rectangle (13, 13);
  % Línea frontal: centro en y = +12 (de +11 a +13)
  \fill[linea] (-13, 11) rectangle (13, 13);
  % Línea trasera: centro en y = -12 (de -13 a -11)
  \fill[linea] (-13, -13) rectangle (13, -11);

  % --- Ejes de referencia (centro de celda = centro del eje de ruedas) ---
  \draw[thin, gristexto!60, dash pattern=on 4pt off 2pt on 1pt off 2pt]
    (0, -14) -- (0, 18.2);
  \draw[thin, gristexto!60, dash pattern=on 4pt off 2pt on 1pt off 2pt]
    (-14.5, 0) -- (14.5, 0);
  \node[font=\scriptsize, gristexto, anchor=south west] at (0.2, 0.2) {Eje $(0,0)$};

  % --- Ruedas (ancho total 15 cm: x de -7.5 a +7.5; diámetro 4.3 cm, ancho 1.9 cm) ---
  \filldraw[rueda] (-7.5, -2.15) rectangle (-5.6, 2.15);
  \filldraw[rueda] (5.6, -2.15) rectangle (7.5, 2.15);
  % Eje mecánico
  \draw[thick, black!70] (-5.6, 0) -- (5.6, 0);

  % --- Chasis (18 cm de largo x 10 cm de ancho: x in [-5,5], y in [-4,14]) ---
  \filldraw[chasis] (-5, -4) rectangle (5, 14);
  \node[font=\footnotesize\bfseries, azuldoc, align=center] at (0, 5.2)
    {Chasis 2WD\\$18 \times 10\text{ cm}$};

  % --- Soporte frontal hasta y = 16 cm (largo total = 20 cm) ---
  \draw[azuldoc, line width=1pt, fill=azulfondo!70] (-1.2, 14) rectangle (1.2, 16);

  % --- Sensores TCRT5000 (3 unidades) ---
  % Laterales: 16 cm entre puntas (x = -8 y x = +8)
  \draw[thick, azuldoc] (-5, 7.5) -- (-8, 7.5);
  \draw[thick, azuldoc] (5, 7.5) -- (8, 7.5);
  \filldraw[sensor] (-8.4, 7.0) rectangle (-7.6, 8.0);
  \filldraw[sensor] (7.6, 7.0) rectangle (8.4, 8.0);
  % Sensor frontal: a 16 cm del eje (y = +16)
  \filldraw[sensor] (-0.6, 15.6) rectangle (0.6, 16.4);

  % Etiquetas de sensores (dentro del espacio blanco, sin tocar líneas negras)
  \node[font=\scriptsize\bfseries, rojoborde, anchor=south, fill=white, inner sep=1pt]
    at (-8.0, 8.3) {TCRT Izq.};
  \node[font=\scriptsize\bfseries, rojoborde, anchor=south, fill=white, inner sep=1pt]
    at (8.0, 8.3) {TCRT Der.};
  \node[font=\scriptsize\bfseries, rojoborde, anchor=west, fill=white, inner sep=1pt]
    at (1.3, 16.3) {TCRT Frontal};

  % --- Flecha de sentido de avance ---
  \draw[->, very thick, unalverde] (0, 8.5) -- (0, 12.2);
  \node[font=\scriptsize\bfseries, unalverde, anchor=west] at (0.3, 10.2) {Avance};

  % ==========================================================================
  % ACOTACIONES (sin solapamientos)
  % ==========================================================================

  % 1. Celda 24 cm (centro a centro) - parte inferior externa
  \draw[ext] (-12, -13) -- (-12, -17.5);
  \draw[ext] (12, -13) -- (12, -17.5);
  \draw[cota] (-12, -16.8) -- (12, -16.8)
    node[midway, fill=white, inner sep=2pt, font=\footnotesize\bfseries]
    {Celda: \SI{24}{cm} (centro a centro)};

  % 2. Blanco interior 22 cm - parte inferior interna
  \draw[ext] (-11, -11) -- (-11, -14.8);
  \draw[ext] (11, -11) -- (11, -14.8);
  \draw[cota] (-11, -14.2) -- (11, -14.2)
    node[midway, fill=white, inner sep=2pt, font=\scriptsize]
    {Blanco interior: \SI{22}{cm}};

  % 3. Ancho con ruedas: 15 cm (abajo del chasis)
  \draw[ext] (-7.5, -2.15) -- (-7.5, -7.2);
  \draw[ext] (7.5, -2.15) -- (7.5, -7.2);
  \draw[cota] (-7.5, -6.5) -- (7.5, -6.5)
    node[midway, fill=white, inner sep=2pt, font=\scriptsize\bfseries]
    {Ancho con ruedas: \SI{15}{cm}};

  % 4. Distancia entre sensores laterales: 16 cm (por encima del chasis interno)
  \draw[ext] (-8, 7.0) -- (-8, 3.2);
  \draw[ext] (8, 7.0) -- (8, 3.2);
  \draw[cota] (-8, 3.6) -- (8, 3.6)
    node[midway, fill=azulfondo, inner sep=2pt, font=\scriptsize\bfseries]
    {\SI{16}{cm} entre sensores laterales};

  % 5. Tolerancia lateral sensores: 3 cm a la izquierda (x = -11 a -8)
  \draw[ext] (-11, 7.5) -- (-11, 4.8);
  \draw[ext] (-8, 7.0) -- (-8, 4.8);
  \draw[cota, rojoborde] (-11, 5.2) -- (-8, 5.2);
  \node[font=\scriptsize\bfseries, rojoborde, anchor=south, fill=white, inner sep=1pt]
    at (-9.5, 5.45) {\SI{3}{cm}};

  % 6. Tolerancia lateral ruedas: 3.5 cm a la derecha (x = 7.5 a 11)
  \draw[ext] (7.5, -1.5) -- (7.5, -4.6);
  \draw[ext] (11, -1.5) -- (11, -4.6);
  \draw[cota, unalverde] (7.5, -4.2) -- (11, -4.2);
  \node[font=\scriptsize\bfseries, unalverde, anchor=west, fill=white, inner sep=1pt]
    at (13.4, -4.2) {Holgura rueda: \SI{3.5}{cm}};
  \draw[thin, unalverde] (11.1, -4.2) -- (13.3, -4.2);

  % 7. Grosor de línea: 2 cm (borde derecho inferior, sin cruzar Media celda)
  \draw[cota] (11, -8.5) -- (13, -8.5);
  \node[font=\scriptsize, anchor=west, fill=white, inner sep=1pt]
    at (13.4, -8.5) {Línea: \SI{2}{cm}};
  \draw[thin, azuldoc] (13.3, -8.5) -- (13.05, -8.5);

  % 8. Distancia eje -> sensor frontal: 16 cm (lado izquierdo externo)
  \draw[ext] (0, 0) -- (-15.8, 0);
  \draw[ext] (-0.6, 16) -- (-15.8, 16);
  \draw[cota] (-15.2, 0) -- (-15.2, 16)
    node[midway, rotate=90, fill=white, inner sep=2pt, font=\scriptsize\bfseries]
    {Eje a sensor frontal: \SI{16}{cm}};

  % 9. Largo total del robot: 20 cm (lado izquierdo más externo)
  \draw[ext] (-5, -4) -- (-18.5, -4);
  \draw[ext] (-0.6, 16) -- (-18.5, 16);
  \draw[cota] (-17.8, -4) -- (-17.8, 16)
    node[midway, rotate=90, fill=white, inner sep=2pt, font=\scriptsize\bfseries]
    {Largo total: \SI{20}{cm}};

  % 10. Saliente de 4 cm del sensor frontal respecto al centro de línea (y = 12 a 16)
  \draw[ext] (0, 12) -- (15.5, 12);
  \draw[ext] (0.6, 16) -- (15.5, 16);
  \draw[cota, rojoborde] (14.8, 12) -- (14.8, 16)
    node[midway, anchor=west, fill=white, inner sep=2pt, font=\scriptsize\bfseries, align=left]
    {Saliente: \SI{4}{cm}\\(sobre centro de línea)};

  % 11. Media celda 12 cm (lado derecho, y = 0 a 12)
  \draw[ext] (7.5, 0) -- (15.5, 0);
  \draw[cota] (14.8, 0) -- (14.8, 12)
    node[midway, rotate=90, fill=white, inner sep=2pt, font=\scriptsize]
    {Media celda: \SI{12}{cm}};

\end{tikzpicture}

  \caption{Vista superior acotada del robot diferencial centrado en una celda de \SI{24}{\centi\meter} con líneas de \SI{2}{\centi\meter}. Se muestran las holguras laterales de sensores (\SI{3}{\centi\meter}) y ruedas (\SI{3.5}{\centi\meter}), así como el saliente frontal de \SI{4}{\centi\meter} respecto al centro de la línea.}
  \label{fig:geometria-carro}
\end{figure}

\begin{decision}{DEC-02: Tamaño de celda y grosor de línea}
Fijamos una celda de \SI{24}{\centi\meter} entre centros de líneas, trazo de \SI{2}{\centi\meter} y margen exterior blanco de \SI{5}{\centi\meter}. Así quedan $24 - 2 = \SI{22}{\centi\meter}$ de blanco entre paredes. Como los TCRT5000 laterales están separados \SI{16}{\centi\meter}, el carro tiene $\frac{22 - 16}{2} = \SI{3}{\centi\meter}$ de holgura por lado antes de pisar una línea, y las ruedas (\SI{15}{\centi\meter} de ancho total) cuentan con $\SI{3.5}{\centi\meter}$ por lado.

El margen de \SI{5}{\centi\meter} alrededor del borde no es estético: el sensor frontal sobresale \SI{16}{\centi\meter} del eje y media celda mide \SI{12}{\centi\meter}, así que al quedar centrado el sensor pasa \SI{4}{\centi\meter} más allá de la línea (y la punta del chasis unos \SI{5}{\centi\meter}). Sin ese margen blanco el sensor se saldría del papel en las paredes del perímetro~\cite{arquitectura2026doc}.

Para la primera prueba del 26/10/2026 sobre \archivo{maps/simulados/prueba\_4x4.txt} tuvimos que encoger la pista a $91 \times \SI{91}{\centi\meter}$ (celda de \SI{22}{\centi\meter}, línea de \SI{2}{\centi\meter}, margen de \SI{0.5}{\centi\meter}) para que cupiera en un rollo de plotter de \SI{91}{\centi\meter}. Eso dejó apenas \SI{2}{\centi\meter} de margen por sensor y \SI{2.5}{\centi\meter} por rueda~\cite{handoff2026doc}, lo que confirmó que \SI{24}{\centi\meter} es el tamaño cómodo para la demo.
\end{decision}

% ============================================================================
\subsection{Respuesta infrarroja y elección del material del laberinto (\cmd{DEC-03})}
\label{subsec:ir-materiales}

Los tres TCRT5000~\cite{vishay2012tcrt5000} entran al ADC1 de 12~bits (\numrange{0}{4095}). En el firmware normalizamos cada lectura a un entero entre $0$ y $1000$:
\begin{equation}
  S_{\text{norm}} = \left\lfloor \frac{S_{\text{crudo}} \cdot 1000}{4095} \right\rfloor,
  \label{eq:adc-norm}
\end{equation}
donde los números bajos corresponden a blanco (mucha reflexión IR) y los altos a negro mate o al sensor en el aire. La \Cref{tab:lecturas-ir} reúne las lecturas reales que medimos el 23/10 y el 26/10/2026.

\begin{table}[H]
  \centering
  \caption{Lecturas normalizadas (\numrange{0}{1000}) de los sensores TCRT5000 sobre distintas superficies.}
  \label{tab:lecturas-ir}
  \renewcommand{\arraystretch}{1.15}
  \small
  \begin{tabularx}{\textwidth}{@{}p{4.4cm} c c X@{}}
    \toprule
    \textbf{Superficie / Condición} & \textbf{Rango normalizado} & \textbf{Fecha} & \textbf{Interpretación del sensor (\cmd{umbral = 200})} \\
    \midrule
    Papel blanco (impreso / base) & \numrange{35}{57} / \numrange{24}{85} & 23/10 -- 26/10 & Blanco (alta reflexión IR) \\
    Cinta aislante negra mate & \numrange{250}{860} & 23/10/2026 & Negro (absorbe IR; válida para detección) \\
    Cinta aislante sobre impreso & \numrange{400}{690} & 26/10/2026 & Negro claro (margen $> \num{300}$ sobre blanco) \\
    Aire (robot levantado) & \num{1000} & 23/10/2026 & Saturación por ausencia de retorno IR \\
    Tinta negra de plotter & $\approx$~\numrange{24}{85} & 26/10/2026 & \textbf{No sirve:} lectura igual al papel blanco \\
    \emph{Mousepad} negro brillante & $< \num{200}$ & 26/10/2026 & \textbf{No sirve:} refleja el IR como si fuera blanco \\
    \bottomrule
  \end{tabularx}
\end{table}

Con estos valores pusimos el umbral inicial en \cmd{umbral = 200} en \archivo{parametros.h}, a mitad de camino entre el blanco más alto (\num{85}) y la cinta más baja (\numrange{250}{400}).

\begin{hallazgo}[Hallazgo del 26/10/2026: la tinta del plotter es invisible al IR]
  Cuando pusimos el carro sobre el primer laberinto recién salido del plotter, los tres TCRT5000 pasaron por encima de las líneas negras sin enterarse: marcaban entre \num{24} y \num{85}, igual que en el papel blanco, incluso bajando el sensor a pocos milímetros. Esa tinta refleja en la banda de \SI{950}{\nano\meter} igual que el papel. Algo parecido nos pasó con un \emph{mousepad} negro brillante~\cite{repo2026agentelaberinto,arquitectura2026doc}.
\end{hallazgo}

\begin{decision}{DEC-03: Material físico del laberinto}
  El plano del plotter salido de \archivo{pc/imprimir.py} sirve como guía de medidas, pero las paredes se forran con cinta aislante negra mate ($\approx\SI{1.9}{\centi\meter}$ centrada sobre el trazo de \SI{2}{\centi\meter}, sin arrugas y prolongada $\sim\SI{1}{\centi\meter}$ en las esquinas), o bien se arma todo con tiras de cartulina negra mate probadas antes con el sensor ($> 400$)~\cite{arquitectura2026doc}.
  \par\smallskip
  Queda un punto abierto para exteriores (\cmd{HW-04}): hoy los sensores cuelgan a $\sim\SI{1.5}{\centi\meter}$ del piso, una altura en la que el sol directo puede tapar el contraste. El plan antes de la demo al aire libre es bajarlos a \SIrange{5}{8}{\milli\meter} y ponerles una faldita de cartón alrededor~\cite{arquitectura2026doc}.
\end{decision}

% ============================================================================
\subsection{Alimentación y caídas de voltaje}
\label{subsec:alimentacion}

Separamos la energía en dos circuitos independientes con tierra común (GND), porque los picos inductivos de los motores al arrancar y frenar reiniciaban la placa si compartían fuente. La lógica (ESP32-S3, MPU6050 y TCRT5000 a \SI{3.3}{\volt}, más la lógica de \SI{5}{\volt} del L298N) se alimenta con una batería de \SI{9}{\volt} en el jack DC del shield (\SIrange{6.5}{9}{\volt}). Hay una regla que respetamos siempre al pie de la letra: nunca conectar el cable USB-C y la batería de \SI{9}{\volt} al tiempo~\cite{handoff2026doc}. Para los motores (borne \SI{12}{\volt}/VMS del L298N), la configuración de \texttt{main} usa cuatro pilas AA en serie (\SI{6}{\volt}); en la rama \archivo{origin/firmware/volver} migramos a un pack 2S de celdas Li-Ion 18650 (\SIrange{7.8}{8.4}{\volt} con interruptor general) para recuperar el torque perdido en el L298N.

La batería de \SI{9}{\volt} nos dio un susto temprano. El 23/10/2026, probando \archivo{prueba\_alimentacion\_esp32s3.ino}, una batería que medía \SI{7.67}{\volt} en reposo metió al ESP32 en un bucle infinito de reinicios (\cmd{ESP\_RST\_BROWNOUT}, con el LED verde TX del CH343 parpadeando sin parar) apenas encendía el Wi-Fi. Lo resolvimos bajando la potencia Wi-Fi a \SI{8.5}{dBm} (\cmd{WiFi.setTxPower(WIFI\_POWER\_8\_5dBm)}) y fijando la regla práctica de cambiar la batería de \SI{9}{\volt} en cuanto baje de \SI{8.5}{\volt} (nueva marca cerca de \SI{9.5}{\volt})~\cite{arquitectura2026doc}.

En la rama \archivo{origin/firmware/volver} añadimos además el divisor resistivo en GPIO~1 para leer el pack 2S cada \SI{5}{\milli\second} durante la marcha. El 04/11/2026, con el pack en \SI{8.0}{\volt} en reposo, medimos caídas de $\approx\SI{0.6}{\volt}$ en los giros y $\approx\SI{0.36}{\volt}$ al avanzar. Filtrando esa lectura con $\tau \approx \SI{1}{\second}$ y escalando el PWM por $(v_{\text{ref}} - v_{\text{caida}}) / (V_{\text{bat}} - v_{\text{caida}})$, evitamos tener que recalibrar \cmd{t\_celda} cada media hora conforme se descarga la batería~\cite{repo2026agentelaberinto}.

% ============================================================================
\subsection{Esquema de conexiones y mapa de pines GPIO}
\label{subsec:esquema-pines}

La \Cref{fig:cableado} resume todas las conexiones entre la placa, el L298N, los sensores, el giroscopio y el divisor de batería.

\begin{figure}[H]
  \centering
  % ============================================================================
% Figura: Diagrama eléctrico y de conexiones (ESP32-S3, L298N, sensores y potencia)
% ============================================================================
\begin{circuitikz}[
  scale=0.61, transform shape,
  >=Stealth,
  bloque/.style={draw=azuldoc, fill=azulfondo, rounded corners=3pt, line width=1pt},
  potencia/.style={draw=rojoborde, fill=rojofondo, rounded corners=3pt, line width=1pt},
  sensor/.style={draw=unalverde, fill=verdefondo, rounded corners=3pt, line width=1pt},
  alerta/.style={draw=ambarborde, fill=ambarfondo, rounded corners=2pt, line width=0.8pt},
  cable/.style={thick, azuldoc},
  pwr/.style={thick, rojoborde},
  pwr5/.style={thick, orange!85!black},
  pwr3/.style={thick, unalverde},
  gnd/.style={thick, black!80}
]

  % ==========================================================================
  % BLOQUE CENTRAL: ESP32-S3 N16R8 + SHIELD DE EXPANSIÓN
  % ==========================================================================
  \draw[bloque] (-3.2, -5.2) rectangle (3.2, 4.8);
  \node[font=\small\bfseries, azuldoc, align=center] at (0, 4.25)
    {ESP32-S3 N16R8 + Shield};
  \node[font=\scriptsize, gristexto] at (0, 3.65)
    {44 pines $\cdot$ ADC1 $\cdot$ Wi-Fi 8.5 dBm};

  % Pines lado izquierdo (Sensores, IMU, Batería, Alimentación lógica)
  \node[font=\scriptsize, anchor=west] (p3v3) at (-3.1, 2.7) {3.3\,V (G/V/S)};
  \node[font=\scriptsize, anchor=west] (pg4)  at (-3.1, 1.9) {GPIO 4 (ADC1)};
  \node[font=\scriptsize, anchor=west] (pg5)  at (-3.1, 1.3) {GPIO 5 (ADC1)};
  \node[font=\scriptsize, anchor=west] (pg6)  at (-3.1, 0.7) {GPIO 6 (ADC1)};
  \node[font=\scriptsize, anchor=west] (pg9)  at (-3.1, -0.4) {GPIO 9 (SDA)};
  \node[font=\scriptsize, anchor=west] (pg10) at (-3.1, -1.0) {GPIO 10 (SCL)};
  \node[font=\scriptsize, anchor=west] (pg1)  at (-3.1, -2.5) {GPIO 1 (ADC1)};
  \node[font=\scriptsize, anchor=west] (pjack) at (-3.1, -4.1) {Jack DC (6.5--9\,V)};

  % Pines lado derecho (L298N, 5V, LED RGB, GND)
  \node[font=\scriptsize, anchor=east] (pg7)  at (3.1, 2.5) {GPIO 7 (PWM)};
  \node[font=\scriptsize, anchor=east] (pg15) at (3.1, 1.9) {GPIO 15};
  \node[font=\scriptsize, anchor=east] (pg16) at (3.1, 1.3) {GPIO 16};
  \node[font=\scriptsize, anchor=east] (pg8)  at (3.1, 0.3) {GPIO 8 (PWM)};
  \node[font=\scriptsize, anchor=east] (pg17) at (3.1, -0.3) {GPIO 17};
  \node[font=\scriptsize, anchor=east] (pg18) at (3.1, -0.9) {GPIO 18};
  \node[font=\scriptsize, anchor=east] (p5v)  at (3.1, -2.2) {Riel 5\,V};
  \node[font=\scriptsize, anchor=east] (pg48) at (3.1, -3.3) {GPIO 48};
  \node[font=\scriptsize, anchor=east] (pgnd) at (3.1, -4.4) {GND};

  % ==========================================================================
  % COLUMNA IZQUIERDA: SENSORES TCRT5000, MPU6050, DIVISOR BAT, BATERÍA 9V
  % ==========================================================================

  % --- 3x TCRT5000 ---
  \draw[sensor] (-9.6, 0.3) rectangle (-5.2, 3.2);
  \node[font=\footnotesize\bfseries, unalverde, align=center] at (-7.4, 2.75)
    {3$\times$ TCRT5000 (AO)};
  \node[font=\scriptsize, anchor=east] at (-5.3, 1.9) {Izq.};
  \node[font=\scriptsize, anchor=east] at (-5.3, 1.3) {Cen.};
  \node[font=\scriptsize, anchor=east] at (-5.3, 0.7) {Der.};
  \node[font=\scriptsize, gristexto, anchor=west] at (-9.4, 1.6) {VCC: 3.3\,V};
  \node[font=\scriptsize, gristexto, anchor=west] at (-9.4, 0.9) {GND};

  \draw[cable] (-5.2, 1.9) -- (-3.2, 1.9);
  \draw[cable] (-5.2, 1.3) -- (-3.2, 1.3);
  \draw[cable] (-5.2, 0.7) -- (-3.2, 0.7);
  \draw[pwr3] (-3.2, 2.7) -- (-4.5, 2.7) |- (-5.2, 2.5);
  \node[font=\tiny\bfseries, unalverde, fill=white, inner sep=1pt] at (-4.1, 2.9) {3.3\,V};

  % --- MPU6050 (GY-521) ---
  \draw[sensor] (-9.6, -1.6) rectangle (-5.2, -0.05);
  \node[font=\footnotesize\bfseries, unalverde, align=center] at (-7.4, -0.38)
    {MPU6050 (GY-521)};
  \node[font=\scriptsize, gristexto, align=left] at (-7.4, -0.95)
    {I2C 400\,kHz (0x68)\\AD0/XDA/XCL/INT: NC};

  \draw[cable] (-5.2, -0.4) -- (-3.2, -0.4);
  \draw[cable] (-5.2, -1.0) -- (-3.2, -1.0);
  \draw[pwr3] (-4.5, 2.7) -- (-4.5, -0.15) -- (-5.2, -0.15);

  % --- Divisor de batería (rama volver) ---
  \draw[alerta] (-10.8, -3.95) rectangle (-4.1, -1.75);
  \node[font=\scriptsize\bfseries, ambarborde, anchor=west] at (-10.65, -2.0)
    {Divisor de batería [rama volver]};

  % Circuito resistivo + capacitor con CircuiTikZ
  \draw[cable] (-5.7, -2.75) -- (-3.2, -2.75);
  \draw (-9.4, -2.75) node[anchor=east, font=\scriptsize\bfseries, rojoborde] {$V_{\text{bat}}(+)$}
    to[R, l={\scriptsize $6.6\text{ k}\Omega$}, -*] (-6.7, -2.75)
    -- (-5.7, -2.75);
  \draw (-6.7, -2.75) to[R, l_={\scriptsize $2.2\text{ k}\Omega$}] (-6.7, -3.8)
    node[ground]{};
  \draw (-5.7, -2.75) to[C, l={\scriptsize $10\text{ nF}$}] (-5.7, -3.8)
    node[ground]{};

  % --- Batería 9V lógica ---
  \draw[potencia] (-9.6, -5.2) rectangle (-5.2, -4.05);
  \node[font=\footnotesize\bfseries, rojoborde, align=center] at (-7.4, -4.4)
    {Batería 9\,V (Lógica)};
  \node[font=\tiny\bfseries, rojoborde, align=center] at (-7.4, -4.9)
    {NUNCA USB + Batería};
  \draw[pwr] (-5.2, -4.1) -- (-3.2, -4.1);

  % ==========================================================================
  % COLUMNA DERECHA: DRIVER L298N, MOTORES N20-10, PACK DE POTENCIA
  % ==========================================================================

  % --- Driver L298N ---
  \draw[potencia] (5.4, -2.9) rectangle (10.3, 3.2);
  \node[font=\small\bfseries, rojoborde] at (7.85, 2.8) {Driver L298N};
  \node[font=\tiny\bfseries, ambarborde, fill=ambarfondo, draw=ambarborde,
        rounded corners=1.5pt, inner sep=2pt] at (7.85, 2.25)
    {Jumpers ENA, ENB y 5V QUITADOS};

  % Entradas lógicas L298N
  \node[font=\scriptsize, anchor=west] at (5.5, 1.55) {ENA};
  \node[font=\scriptsize, anchor=west] at (5.5, 1.05) {IN1};
  \node[font=\scriptsize, anchor=west] at (5.5, 0.55) {IN2};
  \node[font=\scriptsize, anchor=west] at (5.5, -0.05) {ENB};
  \node[font=\scriptsize, anchor=west] at (5.5, -0.55) {IN3};
  \node[font=\scriptsize, anchor=west] at (5.5, -1.05) {IN4};
  \node[font=\scriptsize\bfseries, orange!85!black, anchor=west] at (5.5, -2.0) {Borne 5V};
  \node[font=\scriptsize\bfseries, anchor=west] at (5.5, -2.55) {GND};

  % Salidas potencia L298N
  \node[font=\scriptsize, anchor=east] at (10.2, 1.15) {OUT1--OUT2};
  \node[font=\scriptsize, anchor=east] at (10.2, -0.45) {OUT3--OUT4};
  \node[font=\scriptsize\bfseries, rojoborde, anchor=east] at (10.2, -2.0) {12V / VMS};

  % Conexiones ESP32 -> L298N
  \draw[cable] (3.2, 2.5) -- (4.3, 2.5) |- (5.4, 1.55);
  \draw[cable] (3.2, 1.9) -- (4.1, 1.9) |- (5.4, 1.05);
  \draw[cable] (3.2, 1.3) -- (3.9, 1.3) |- (5.4, 0.55);
  \draw[cable] (3.2, 0.3) -- (4.3, 0.3) |- (5.4, -0.05);
  \draw[cable] (3.2, -0.3) -- (4.1, -0.3) |- (5.4, -0.55);
  \draw[cable] (3.2, -0.9) -- (3.9, -0.9) |- (5.4, -1.05);
  \draw[pwr5] (3.2, -2.2) -- (4.5, -2.2) |- (5.4, -2.0);
  \node[font=\tiny\bfseries, orange!85!black, fill=white, inner sep=1pt]
    at (4.3, -1.75) {5\,V lógico};

  % --- Motores N20-10 A y B ---
  \draw[bloque] (11.6, 0.45) rectangle (15.4, 1.85);
  \node[font=\footnotesize\bfseries, azuldoc, align=center] at (13.5, 1.15)
    {Motor A (N20-10)\\Llanta Derecha};
  \draw[pwr] (10.3, 1.15) -- (11.6, 1.15);

  \draw[bloque] (11.6, -1.15) rectangle (15.4, 0.25);
  \node[font=\footnotesize\bfseries, azuldoc, align=center] at (13.5, -0.45)
    {Motor B (N20-10)\\Llanta Izquierda};
  \draw[pwr] (10.3, -0.45) -- (11.6, -0.45);

  % --- Pack de baterías de motores ---
  \draw[potencia] (11.2, -3.5) rectangle (15.6, -1.65);
  \node[font=\footnotesize\bfseries, rojoborde, align=center] at (13.4, -2.1)
    {Pack Motores (+ SW)};
  \node[font=\scriptsize, align=center] at (13.4, -2.85)
    {4$\times$AA (6\,V) [main]\\2S 18650 (7.8--8.4\,V) [volver]};
  \draw[pwr] (11.2, -2.0) -- (10.3, -2.0);

  % --- LED RGB WS2812 ---
  \draw[alerta] (5.4, -4.05) rectangle (10.3, -3.15);
  \node[font=\scriptsize\bfseries, ambarborde, align=center] at (7.85, -3.6)
    {LED RGB WS2812\\(puente \texttt{RGB} soldado)};
  \draw[cable] (3.2, -3.3) -- (4.5, -3.3) |- (5.4, -3.6);

  % ==========================================================================
  % BUS DE TIERRA COMÚN (GND)
  % ==========================================================================
  \draw[gnd] (-9.6, -5.75) -- (13.4, -5.75);
  \node[font=\footnotesize\bfseries, fill=white, inner sep=2pt] at (0.0, -5.75)
    {Bus de Tierra Común (GND)};

  % Bajadas a GND común
  \draw[gnd] (-7.4, -5.2) -- (-7.4, -5.75);  % Bat 9V
  \draw[gnd] (3.2, -4.4) -- (3.8, -4.4) -- (3.8, -5.75); % Shield GND
  \draw[gnd] (5.4, -2.55) -- (4.9, -2.55) -- (4.9, -5.75); % L298N GND
  \draw[gnd] (13.4, -3.5) -- (13.4, -5.75); % Pack motores GND
  \draw[gnd] (-9.2, 0.3) -- (-9.2, -0.05); % TCRT GND -> MPU6050
  \draw[gnd] (-9.8, -0.8) -- (-10.4, -0.8) |- (-9.6, -5.75); % Sensores/IMU GND

\end{circuitikz}

  \caption{Diagrama eléctrico y de conexiones del robot. Se muestran los dos dominios de energía con GND común, los \emph{jumpers} retirados en el L298N, los \SI{5}{\volt} lógicos desde el shield y el divisor resistivo en GPIO~1 (rama \archivo{firmware/volver}).}
  \label{fig:cableado}
\end{figure}

La \Cref{tab:pines-gpio} detalla la asignación definida en \archivo{arduino/agente/src/config/pines\_esp32s3.h}. El 23/10/2026 verificamos el sentido de giro con \archivo{prueba\_motores\_esp32s3.ino}: \cmd{IN1=HIGH, IN2=LOW} (Motor~A derecho) e \cmd{IN3=HIGH, IN4=LOW} (Motor~B izquierdo) mueven las dos ruedas hacia adelante sin tener que invertir nada en código.

\begin{table}[H]
  \centering
  \caption{Asignación completa de pines GPIO en el ESP32-S3 N16R8 (\archivo{pines\_esp32s3.h}).}
  \label{tab:pines-gpio}
  \renewcommand{\arraystretch}{1.15}
  \small
  \begin{tabularx}{\textwidth}{@{}c p{3.6cm} p{3.5cm} X@{}}
    \toprule
    \textbf{GPIO} & \textbf{Constante en firmware} & \textbf{Subsistema / Pin destino} & \textbf{Función y notas eléctricas} \\
    \midrule
    1  & \cmd{PIN\_BAT} \badgepr & Divisor \SI{6.6}{\kilo\ohm} / \SI{2.2}{\kilo\ohm} + \SI{10}{\nano\farad} & Lectura ADC1 de $V_{\text{bat}}$ ($V_{\text{pin}} \approx V_{\text{bat}}/4$; \SI{8.4}{\volt} $\rightarrow$ \SI{2.1}{\volt}) \\
    4  & \cmd{PIN\_SENSOR\_IZQ} & TCRT5000 Izquierdo (\cmd{AO}) & ADC1; alimentado a \SI{3.3}{\volt} (filas G/V/S del shield) \\
    5  & \cmd{PIN\_SENSOR\_CEN} & TCRT5000 Frontal (\cmd{AO}) & ADC1; alimentado a \SI{3.3}{\volt} (filas G/V/S del shield) \\
    6  & \cmd{PIN\_SENSOR\_DER} & TCRT5000 Derecho (\cmd{AO}) & ADC1; alimentado a \SI{3.3}{\volt} (filas G/V/S del shield) \\
    7  & \cmd{PIN\_DER\_EN} & L298N \cmd{ENA} & Salida PWM de velocidad, Motor~A (llanta derecha) \\
    8  & \cmd{PIN\_IZQ\_EN} & L298N \cmd{ENB} & Salida PWM de velocidad, Motor~B (llanta izquierda) \\
    9  & \cmd{PIN\_IMU\_SDA} & MPU6050 \cmd{SDA} & Bus I2C a \SI{400}{\kilo\hertz} (\SI{3.3}{\volt}; \cmd{AD0/XDA/XCL/INT} sin conectar) \\
    10 & \cmd{PIN\_IMU\_SCL} & MPU6050 \cmd{SCL} & Reloj I2C a \SI{400}{\kilo\hertz} (\cmd{0x68}) \\
    15 & \cmd{PIN\_DER\_AD} & L298N \cmd{IN1} & Control dirección Motor~A (\cmd{HIGH} = adelante) \\
    16 & \cmd{PIN\_DER\_AT} & L298N \cmd{IN2} & Control dirección Motor~A (\cmd{HIGH} = atrás) \\
    17 & \cmd{PIN\_IZQ\_AD} & L298N \cmd{IN3} & Control dirección Motor~B (\cmd{HIGH} = adelante) \\
    18 & \cmd{PIN\_IZQ\_AT} & L298N \cmd{IN4} & Control dirección Motor~B (\cmd{HIGH} = atrás) \\
    48 & \cmd{PIN\_LED} & LED RGB WS2812 & Estado visual (requiere puente \cmd{RGB} soldado en placa clon) \\
    \bottomrule
  \end{tabularx}
\end{table}

Todas las entradas analógicas (GPIO~1, 4, 5 y 6) van en ADC1 porque en el ESP32-S3 el bloque ADC2 no puede usarse mientras el Wi-Fi está activo~\cite{espressif2024esp32s3}. La \Cref{tab:pines-restringidos} recuerda además qué pines del módulo dejamos libres por restricciones internas del SoC.

\begin{table}[H]
  \centering
  \caption{Pines GPIO restringidos en el módulo ESP32-S3 N16R8 (\archivo{pines\_esp32s3.h}).}
  \label{tab:pines-restringidos}
  \renewcommand{\arraystretch}{1.15}
  \small
  \begin{tabularx}{\textwidth}{@{}l l X@{}}
    \toprule
    \textbf{Pines GPIO} & \textbf{Función reservada} & \textbf{Motivo de restricción} \\
    \midrule
    GPIO 35, 36, 37 & OPI PSRAM (\SI{8}{\mega\byte}) & Bus octal interno de memoria; usarlos bloquea el arranque o corrompe datos \\
    GPIO 0, 3, 45, 46 & Pines de \emph{strapping} & Definen el modo de arranque y voltaje de flash al encender el SoC \\
    GPIO 19, 20 & USB nativo (D$-$ / D$+$) & Reservados para la interfaz USB OTG nativa del ESP32-S3 \\
    GPIO 43, 44 & UART0 (TX / RX) & Conectados al puente USB-UART CH343 para consola serial y flasheo \\
    \bottomrule
  \end{tabularx}
\end{table}

Por último, la \Cref{fig:fotos-hardware} deja reservados los dos espacios donde iremos pegando las fotografías reales del chasis armado y del cableado inferior.

\begin{figure}[H]
  \centering
  \begin{minipage}[t]{0.48\textwidth}
    \placeholderfoto{Vista superior del robot ensamblado}{Fotografía en planta del chasis 2WD con el ESP32-S3 + shield, driver L298N, MPU6050 sobre espuma, pack de baterías y distribución de los tres sensores TCRT5000.}
  \end{minipage}\hfill
  \begin{minipage}[t]{0.48\textwidth}
    \placeholderfoto{Detalle de cableado y sensores inferiores}{Vista inferior/lateral mostrando la altura de montaje de los TCRT5000 ($\sim\SI{1.5}{\centi\meter}$), las conexiones al borne \SI{5}{\volt} del L298N sin \emph{jumpers} y el bus de tierra común.}
  \end{minipage}
  \caption{Espacios reservados para el registro fotográfico del hardware físico.}
  \label{fig:fotos-hardware}
\end{figure}
